\subsection{Universally positive kernels}

In this issue, X is a separated topological space.

Theorem 1. --- Let \(f \in \mathcal{C}(X \times X)\) . The following conditions are equivalent:

\begin{enumerate}
\def\labelenumi{(\roman{enumi})}
\tightlist
\item
  For every compact subset Y of X and every positive measure \(\mu\) on Y, the endomorphism of \(\mathrm{L}^{2}(\mathrm{Y},\mu)\) defined by the kernel \(f|(\mathrm{Y}\times\mathrm{Y})\) (def. 1 of III, p. 29) is positive, in other words, we have
\end{enumerate}

\[
\int_ {\mathrm{Y} \times \mathrm{Y}} \overline {{h (x)}} h (y) f (x, y) d (\mu \otimes \mu) (x, y) \geqslant 0
\]

for all \(h \in \mathcal{L}^{2}(\mathrm{Y}, \mu)\) ;

\begin{enumerate}
\def\labelenumi{(\roman{enumi})}
\setcounter{enumi}{1}
\tightlist
\item
  For every integer \(n \in N\) , every family \((x_{i})_{0 \leqslant i \leqslant n}\) in X and any family \((t_{i})_{0 \leqslant i \leqslant n}\) of complex numbers, we have
\end{enumerate}

\[
\sum_ {i = 0} ^ {n} \sum_ {j = 0} ^ {n} \bar {t} _ {i} t _ {j} f (x _ {i}, x _ {j}) \geqslant 0;
\]

\begin{enumerate}
\def\labelenumi{(\roman{enumi})}
\setcounter{enumi}{2}
\item
  There exists a complex Hilbert space E and a continuous map \(g \colon X \to E\) of total image such that \(f(x, y) = \langle g(x) | g(y) \rangle\) for all \(x\) and \(y\) in \(X\) ;
\item
  There exists a complex Hilbert space E and a continuous map \(g: X \to E\) such that \(f(x, y) = \langle g(x) | g(y) \rangle\) for all x and y in X.
\end{enumerate}

It is clear that (iii) implies (iv), and one sees that (i) implies (ii) by considering the discrete measure \(\mu\) which is the image of the counting measure on \(\{0,\ldots,n\}\) by the mapping \(i\mapsto x_{i}\) and the function h such that

\[
h(x) = \sum_{\substack{0\leqslant i\leqslant n\\ x_{i} = x}}t_{i}.
\]

Let us prove that (iv) implies (i). Suppose there exists a complex Hilbert space E and a continuous map \(g \colon \mathrm{X} \to \mathrm{E}\) such that \(f(x, y) = \langle g(x) \mid g(y) \rangle\) for all \((x, y) \in \mathrm{X} \times \mathrm{X}\) . Let Y be a compact subset of X, \(\mu\) a positive measure on Y and \(h \in \mathcal{L}^{2}(Y, \mu)\) . We have

\[
\begin{array}{r l} & {\int_ {\mathrm{Y} \times \mathrm{Y}} \overline {{h (x)}} h (y) f (x, y) d (\mu \otimes \mu) (x, y)} \\ & {\qquad = \int_ {\mathrm{Y} \times \mathrm{Y}} \overline {{h (x)}} h (y) \langle g (x) | g (y) \rangle d (\mu \otimes \mu) (x, y)} \\ & {\qquad = \left\langle \int_ {\mathrm{Y}} h (x) g (x) d \mu (x) \Big | \int_ {\mathrm{Y}} h (y) g (y) d \mu (y) \right\rangle \geqslant 0} \end{array}
\]

by INT, V, p. 97, § 8, no. 4, prop. 9.

Let us finally prove that (ii) implies (iii). Let \(\widetilde{\mathrm{E}}\) be the space of complex measures with finite support on X. For \(\mu_{1}\) and \(\mu_{2}\) in \(\widetilde{\mathrm{E}}\) we set

\[
\langle \mu_ {1} \mid \mu_ {2} \rangle = \int_ {\mathrm{X} \times \mathrm{X}} f (x, y) (\overline {{\mu}} _ {1} \otimes \mu_ {2}) (x, y).
\]

The sesquilinear form thus defined on \(\widetilde{\mathrm{E}}\) is a positive Hermitian form. Indeed, let \(\mu \in \widetilde{\mathrm{E}}\) ; there exists a finite family \((x_{i})_{0\leqslant i\leqslant n}\) in X and complex numbers \((t_i)_{0\leqslant i\leqslant n}\) such that \(\mu = \sum_{i=0}^{n} t_i \varepsilon_{x_i}\) . We then have

\[
\langle \mu | \mu \rangle = \sum_ {i = 0} ^ {n} \sum_ {j = 0} ^ {n} \bar {t} _ {i} t _ {j} f (x _ {i}, x _ {j}) \geqslant 0
\]

by hypothesis. One defines \(\widetilde{g}\colon\mathbf{X}\to\widetilde{\mathbf{E}}\) by \(\widetilde{g}(x)=\varepsilon_{x}\) . The image of the map \(\widetilde{g}\) generates \(\widetilde{\mathbf{E}}\) . Moreover, for every \((x,y)\in\mathbf{X}\times\mathbf{X}\) , one has on the one hand \(f(x,y)=\langle\widetilde{g}(x)|\widetilde{g}(y)\rangle\) and on the other hand

\[
\| \widetilde {g} (x) - \widetilde {g} (y) \| ^ {2} = f (x, x) + f (y, y) - f (x, y) - f (y, x),
\]

which implies that \(\widetilde{g}\) is continuous since \(f\) is continuous.

Let E be the separated completion Hilbert space of \(\widetilde{E}\) (EVT, V, p. 8, cor. of prop. 4) and \(g\colon X \to E\) the composition of \(\widetilde{g}\) and of the canonical map \(\widetilde{E} \to E\) . Then the map g is continuous, its image is total in E, and one has \(f(x,y) = \langle g(x) | g(y) \rangle\) for all \((x,y) \in X \times X\) .

The method used to prove that (ii) implies (iii) is called the Gelfand-Naimark-Segal construction.

DEFINITION 1. --- One says that a function \(f \in \mathcal{C}(\mathrm{X} \times \mathrm{X})\) is a universally positive kernel on X if it satisfies the equivalent conditions of Theorem 1.

If f is a universally positive kernel on X, a pair \((E,g)\) satisfying condition (iv) of loc. cit. is called a Hilbertian realization of f; if condition (iii) is satisfied, one says that it is a cyclic Hilbertian realization.

We denote \(\mathrm{Noy}_{+}(\mathrm{X})\) the set of universally positive kernels on X.

Let X' be a separated topological space and \(h: X \to X'\) a continuous map. The map \(f \mapsto f \circ (h, h)\) of \(\mathcal{C}(X' \times X')\) in \(\mathcal{C}(X \times X)\) induces by passage to subspaces a map \(\mathrm{Noy}_{+}(X') \to \mathrm{Noy}_{+}(X)\) .

PROPOSITION 1. --- Let f be a universally positive kernel on X. Let \((\mathrm{E}_{1}, g_{1})\) and \((\mathrm{E}_{2}, g_{2})\) be Hilbertian realizations of f. One supposes that \((\mathrm{E}_{1}, g_{1})\) is a cyclic Hilbertian realization. There then exists a unique continuous linear map \(u: E_{1} \to E_{2}\) such that \(g_{2} = u \circ g_{1}\) . This map is isometric. If \((\mathrm{E}_{2}, g_{2})\) is also cyclic, then u is an isomorphism.

The uniqueness of \(u\) follows from the fact that the image of \(g_{1}\) is total in \(\mathrm{E}_{1}\) . Let \(\mathrm{F} = \mathbf{C}^{(\mathrm{X})}\) and let \((e_x)_{x\in \mathrm{X}}\) the canonical basis of \(\mathrm{F}\) . For \(j = 1\) and \(j = 2\) , denote \(u_{j}\) the linear map from \(\mathrm{F}\) in \(\mathrm{E}_{j}\) determined by \(u_{j}(e_{x}) = g_{j}(x)\) , and denote by \(\mathrm{F}_{j}\) its image; the space \(\mathrm{F}_{1}\) is dense in \(\mathrm{E}_{1}\) . Let \(t = \sum t_{x}e_{x}\) be an element of \(\mathrm{F}\) . We have

\[
\begin{array}{c} \| u _ {1} (t) \| ^ {2} = \sum_ {x, y} \bar {t} _ {x} t _ {y} \langle g _ {1} (x)   |   g _ {1} (y) \rangle = \sum_ {x, y} \bar {t} _ {x} t _ {y} f (x, y) \\ = \sum_ {x, y} \bar {t} _ {x} t _ {y} \langle g _ {2} (x)   |   g _ {2} (y) \rangle = \| u _ {2} (t) \| ^ {2}. \end{array}
\]

Consequently, there exists a linear isometric map v from \(F_{1}\) in \(F_{2}\) such that \(u_{2}=v\circ u_{1}\) , and in particular \(g_{2}(x)=v(g_{1}(x))\) for all \(x\in X\) . Since the image of \(g_{1}\) is total in \(E_{1}\) , this map extends to an isometric linear map u from \(E_{1}\) in \(E_{2}\) such that \(g_{2}=u\circ g_{1}\) .

By Lemma 8 of I, p. 107, the image of \(u\) is closed in \(\mathrm{E}_2\) . If \((\mathrm{E}_2,g_2)\) is a cyclic Hilbert realization, the image of \(u\) is also dense in \(\mathrm{E}_2\) and \(u\) is therefore an isomorphism.

PROPOSITION 2. --- The set Noy \(_{+}\) (X) is a self-adjoint cone in the space \(\mathcal{C}(\text{X} \times \text{X})\) ; it is stable under product and it is closed when one equips \(\mathcal{C}(\text{X} \times \text{X})\) with the topology of simple convergence.

It is elementary that if \(f \in \mathrm{Noy}_{+}(\mathrm{X})\) , then \(tf \in \mathrm{Noy}_{+}(\mathrm{X})\) for every real number \(t \geqslant 0\) , and \(\overline{f} \in \mathrm{Noy}_{+}(\mathrm{X})\) .

If \((\mathrm{E}_{1}, g_{1})\) and \((\mathrm{E}_{2}, g_{2})\) are Hilbertian realizations of universally positive kernels \(f_{1}\) and \(f_{2}\) on X, then the pair \((\mathrm{E}_{1} \oplus \mathrm{E}_{2}, g_{1} + g_{2})\) (resp. the pair \((\mathrm{E}_{1} \widehat{\otimes}_{2} \mathrm{E}_{2}, g_{1} \otimes g_{2})\) ) is a Hilbertian realization of \(f_{1} + f_{2}\) (resp. of \(f_{1}f_{2}\) ); they are therefore universally positive kernels.

The characterization (ii) of \(\mathrm{Noy}_{+}(\mathrm{X})\) (theorem 1) implies that this set is closed in \(\mathcal{C}(\mathrm{X}\times\mathrm{X})\) endowed with the topology of simple convergence.

